\section{Hamiltonian ODE infrastructure (shared by tracks I, N, P)}\label{bpN:p3}

This part is the single ODE layer of the three blueprints I, N, P. It has two levels. The \emph{general layer}
(\S\ref{bpN:p3:s-general}--\S\ref{bpN:p3:s-exit}) is about an autonomous vector field $F$ of class $C^1$ on an open
subset $W$ of a finite-dimensional real normed space $E$: local existence, uniqueness on \emph{every} interval,
maximal solutions, continuation (compact sets, and closed sets with uniform bounds), first exit from an open set.
The \emph{Hamiltonian layer} (\S\ref{bpN:p3:s-ham}) adds $F=X_H$: smoothness of $X_H$, energy conservation,
translation in the angles. \S\ref{bpN:p3:s-frozen} shows that the frozen solution predicates of the three tracks
(\Lean{IsSolutionOn}, \Lean{IsHamiltonianSolutionOn}, \Lean{IsRestrictedSolutionOn}, \Lean{IsNewtonSolution},
\Lean{IsFixedSolution}) are instances of the general predicate, and \S\ref{bpN:p3:s-N} gives the six track-N
results of \texttt{PLAN.md}. Tracks I and P cite the general results as \Nref{bpN:p3:gen-\dots}.

\subsection{Conventions and external facts}\label{bpN:p3:s-conv}

\begin{itemize}
\item $E$ is a finite-dimensional real normed space (hence complete and proper: closed bounded sets are compact),
  $F:E\to E$ a \emph{total} function, $W\subset E$ open. In the applications $E=\R^n\times\R^n$ with the max norm
  $\abs{(\theta,I)}=\max(\abs\theta,\abs I)$ (the product norm of the Lean text, equal to the max norm on the $2n$
  coordinates), or $E=(\R^2)^N\times(\R^2)^N$ for the Newtonian problems.
\item $F$ is \emph{$C^1$ on $W$} means $\texttt{ContDiffOn}\ \R\ 1\ F\ W$; as $W$ is open this gives $F\in C^1$ at each
  point of $W$, and no information off $W$. Off $W$ the values of $F$ never enter (Lean totality, below).
\item An \emph{interval} is an order-connected subset of $\R$ (Lean \Lean{OrdConnected}); open, closed, half-open,
  bounded or not, possibly a single point. $\overline B(x,\delta)$ is the closed ball.
\item All derivatives $\dot z(t)$ in solution predicates are two-sided (Lean \Lean{HasDerivAt}) at every $t\in J$,
  including endpoints of $J$, exactly as in the frozen predicates (``$z$ is defined on $\R$'').
\end{itemize}

\paragraph{External facts.} Mathlib-level statements, with the declarations that exist (checked in
\texttt{Mathlib/Analysis/ODE/}); what is \emph{not} in Mathlib is listed after them.
{\raggedright
\begin{enumerate}[label=(E\arabic*),leftmargin=2.2em]
\item\label{bpN:p3:ext-pl} \emph{Picard--Lindel\"of, quantitative form.} Let $F:E\to E$, $x\in E$, $a,M,L\ge0$ with
  $\norm{F(y)}\le M$ and $F$ $L$-Lipschitz on $\overline B(x,a)$, $T\ge0$ with $MT\le a$. Then there is $\alpha:\R\to E$
  with $\alpha(t_0)=x$ and $\alpha$ has derivative $F(\alpha(t))$ within $[t_0-T,t_0+T]$ at each $t$ of that interval.
  (Mathlib: \Lean{ODE.IsPicardLindelof} with $r=0$, time-independent field, and
  \Lean{IsPicardLindelof.exists\_eq\_forall\_mem\_Icc\_hasDerivWithinAt\_0}; the time constraint is the field
  \Lean{mul\_max\_le}: $M\max(t_{\max}-t_0,t_0-t_{\min})\le a-r$. The contraction on iterates means no smallness of
  $LT$ is needed.) Exists.
\item\label{bpN:p3:ext-unique} \emph{Uniqueness with a Lipschitz constant on a set.} If $S\subset E$ and $F$ is
  $L$-Lipschitz on $S$, $f,g:[a,b]\to E$ are continuous, have right derivative $F(f(t))$, $F(g(t))$ on $[a,b)$, take
  values in $S$ and $f(a)=g(a)$, then $f=g$ on $[a,b]$; and the mirror statement with $f(b)=g(b)$ and left
  derivatives. (Mathlib: \Lean{ODE\_solution\_unique\_of\_mem\_Icc\_right}, \Lean{ODE\_solution\_unique\_of\_mem\_Icc\_left}.)
  Exists.
\item\label{bpN:p3:ext-lip-local} \emph{$C^1$ implies locally Lipschitz.} If $F$ is $C^1$ at $p$ there are $L$ and a
  neighbourhood on which $F$ is $L$-Lipschitz (Mathlib: \Lean{ContDiffAt.exists\_lipschitzOnWith}). Exists.
\item\label{bpN:p3:ext-mvt} \emph{Mean value inequality}: \Kref{bp:p1:ext-mvt}; for a differentiable map on a convex set
  with $\norm{DF}\le L$, $F$ is $L$-Lipschitz there; for a function of one variable with derivative bounded by $M$ on an
  interval, $\norm{\alpha(u)-\alpha(s)}\le M\abs{u-s}$. Exists.
\item\label{bpN:p3:ext-compact} \emph{Topology.} Heine--Borel in $E$ (\Lean{ProperSpace}); continuous images of compact
  sets are compact; sequential compactness; $x\mapsto\operatorname{dist}(x,A)$ is continuous and positive on $E\setminus A$
  for $A$ closed (\Lean{Metric.infDist}). Exists.
\item\label{bpN:p3:ext-calc} \emph{Calculus.} Chain rule \Lean{HasFDerivAt.comp\_hasDerivAt}; a function continuous on
  $[s,t]$ with right derivative $0$ on $[s,t)$ is constant (\Lean{constant\_of\_has\_deriv\_right\_zero}); two one-sided
  derivatives agree $\Rightarrow$ two-sided one (\Lean{HasDerivWithinAt.union} with $(-\infty,b]\cup[b,\infty)=\R$ and
  \Lean{hasDerivWithinAt\_univ}); \Lean{ContDiffOn.fderiv\_of\_isOpen}: if $H$ is $C^{k+1}$ on an open set then $DH$ is
  $C^k$ there. Exists.
\end{enumerate}}
\textbf{Not in Mathlib} (checked: no maximal-solution, continuation, first-exit or uniform-existence-time statement
in \texttt{Analysis/ODE}): everything of \S\ref{bpN:p3:s-general}--\S\ref{bpN:p3:s-exit} beyond (E1)--(E3), notably
Lemmas~\ref{bpN:p3:gen-lip}, \ref{bpN:p3:gen-short}, Proposition~\ref{bpN:p3:gen-uniform-exist},
Theorems~\ref{bpN:p3:gen-uniqueness}, \ref{bpN:p3:gen-maximal}, \ref{bpN:p3:gen-continuation}, and the Hamiltonian
layer. They must be proved from the blueprint below.

\begin{remark}[Lean totality]\label{bpN:p3:totality}
\Lean{KAM.hamiltonianVectorField}\,$H$ uses \Lean{fderiv}, which is $0$ at a point of non-differentiability, and the
Newtonian fields use $1/\ell^3$ with $\ell=\texttt{planeLength}$ or \Lean{radius}. In every use below these are guarded:
$X_H$ occurs only on an open set where $H$ is $C^2$ (or inside the frozen predicates, which assert differentiability
at every visited point), and $\ell>0$ at every visited point because the frozen predicates require it.
The general layer never uses $F$ off $W$.
\end{remark}

\subsection{The general solution predicate; Lipschitz data}\label{bpN:p3:s-general}

\begin{definition}[Solutions in $W$]\label{bpN:p3:gen-sol}
For $F:E\to E$, $W\subset E$, $J\subset\R$ and $z:\R\to E$: $\texttt{IsSolIn}\,F\,W\,J\,z$ holds if for every $t\in J$,
$z(t)\in W$ and $z$ has derivative $F(z(t))$ at $t$ (two-sided). Suggested Lean name
\Lean{KAM.ODE.IsSolutionIn}; it is \emph{not} a frozen declaration.
\source{Standard; shape dictated by the frozen predicates of Definition~\ref{bpN:p0:sol} and \S\ref{bpN:p3:s-frozen}.}
\end{definition}

\noindent Elementary consequences (each a one-line Lean lemma, \ref{bpN:p3:ext-calc}):
(i) restriction: $\texttt{IsSolIn}\,F\,W\,J\,z$ and $J'\subset J$ give $\texttt{IsSolIn}\,F\,W\,J'\,z$;
(ii) $z$ is continuous on $J$ (differentiable at each point of $J$);
(iii) \emph{local nature}: if $J'$ is a neighbourhood of each of its points and $z=y$ on it, then
$\texttt{IsSolIn}\,F\,W\,J'\,y\Leftrightarrow\texttt{IsSolIn}\,F\,W\,J'\,z$ for open $J'$.

\begin{lemma}[Symmetries]\label{bpN:p3:gen-symmetry}
Let $\texttt{IsSolIn}\,F\,W\,J\,z$.
\begin{enumerate}[label=(\alph*)]
\item (time shift) For $c\in\R$, $t\mapsto z(t+c)$ is a solution on $J-c$.
\item (time reversal) $t\mapsto z(-t)$ is a solution of $-F$ on $-J$, and $-F$ is $C^1$ on $W$ if $F$ is.
\item (space translation) If $v\in E$ satisfies $W+v=W$ and $F(x+v)=F(x)$ for $x\in W$, then $t\mapsto z(t)+v$ is a
  solution of $F$ in $W$ on $J$.
\end{enumerate}
\source{Standard.}
\end{lemma}
\begin{proof}
\proofstep{Chain rule}\label{step:p3:sym} (a) and (b) are the chain rule for $t\mapsto t+c$, $t\mapsto-t$ (\ref{bpN:p3:ext-calc});
(c): $z+v$ has the same derivative $F(z(t))=F(z(t)+v)$, and $z(t)+v\in W+v=W$.
\end{proof}

\begin{lemma}[Bounds and Lipschitz constants on compact sets]\label{bpN:p3:gen-lip}
Let $F$ be $C^1$ on the open set $W$ and $K\subset W$ compact. Then $F$ is bounded on $K$ and there is $L\ge0$ with
$\norm{F(x)-F(y)}\le L\norm{x-y}$ for $x,y\in K$.
\source{Standard.}
\end{lemma}
\begin{proof}
\proofstep{Bounded}\label{step:p3:lip-bdd} $F$ is continuous on $W\supset K$, so $\norm F\le M$ on the compact $K$ (\ref{bpN:p3:ext-compact}).

\proofstep{Contradiction hypothesis}\label{step:p3:lip-contra}
If no $L$ works, for every $k\in\N$ there are $x_k\ne y_k$ in $K$ with $\norm{F(x_k)-F(y_k)}>k\norm{x_k-y_k}$; then
$\norm{x_k-y_k}<2M/k\to0$.

\proofstep{Compactness}\label{step:p3:lip-compact}
A subsequence $x_{k_j}\to p\in K$, hence also $y_{k_j}\to p$. By \ref{bpN:p3:ext-lip-local} there are $L_p$ and a
neighbourhood $N_p$ of $p$ on which $F$ is $L_p$-Lipschitz. For large $j$ both points lie in $N_p$ and $k_j>L_p$, which
contradicts Step~\ref{step:p3:lip-contra}.
\end{proof}

\begin{definition}[Uniform data]\label{bpN:p3:gen-uniform}
Let $F$ be $C^1$ on the open set $W$. A set $\Phi\subset E$ has \emph{uniform data $(\delta,M,L)$}, with $\delta>0$,
$M,L\ge0$, if for every $x\in\Phi$: $\overline B(x,\delta)\subset W$, $\norm{F}\le M$ on $\overline B(x,\delta)$ and
$F$ is $L$-Lipschitz on $\overline B(x,\delta)$. (Balls are convex, so by \ref{bpN:p3:ext-mvt} a bound
$\norm{DF}\le L$ on $\overline B(x,\delta)\subset W$ gives the Lipschitz clause.)
\end{definition}

\begin{lemma}[Compact sets have uniform data]\label{bpN:p3:gen-uniform-compact}
Let $F$ be $C^1$ on the open set $W$ and $K\subset W$ compact. Then $K$ has uniform data $(\delta,M,L)$ for some
$\delta>0$, and $K_\delta=\{y:\operatorname{dist}(y,K)\le\delta\}\subset W$.
\end{lemma}
\begin{proof}
\proofstep{Choice of $\delta$}\label{step:p3:uc-delta} If $W=E$ put $\delta=1$. Otherwise $x\mapsto\operatorname{dist}(x,E\setminus W)$ is continuous and
positive on $K$, so has a minimum $2\delta>0$ (\ref{bpN:p3:ext-compact}); then $K_\delta\subset W$.
$K_\delta$ is closed and bounded, hence compact.

\proofstep{Data}\label{step:p3:uc-data} Apply Lemma~\ref{bpN:p3:gen-lip} to $K_\delta$: $\norm F\le M$ and $F$ is $L$-Lipschitz on $K_\delta\supset\overline B(x,\delta)$ for $x\in K$.
\end{proof}

\begin{lemma}[Uniform data for angle-periodic fields]\label{bpN:p3:gen-uniform-periodic}
Let $E=\R^n\times\R^n$ (max norm), $U\subset\R^n$ open, $W=\R^n\times U$, $F$ $C^1$ on $W$ with
$F(\theta+k,I)=F(\theta,I)$ for $(\theta,I)\in W$, $k\in\Z^n$. If $K\subset U$ is compact, then $\Phi=\R^n\times K$ has
uniform data $(\delta,M,L)$ for some $\delta\in(0,1]$.
\end{lemma}
\begin{proof}
\proofstep{A compact cell}\label{step:p3:up-cell} Choose $\delta\in(0,1]$ with $K_\delta=\{I:\operatorname{dist}(I,K)\le\delta\}\subset U$ (as in
Step~\ref{step:p3:uc-delta}). $C=[-1,2]^n\times K_\delta$ is compact and contained in $W$; let $M,L_C$ be the constants of
Lemma~\ref{bpN:p3:gen-lip} for $C$.

\proofstep{Translating a ball into the cell}\label{step:p3:up-translate} Let $x=(\theta,I)\in\Phi$, $k=(\lfloor\theta_i\rfloor)_i\in\Z^n$ and $x'=x-(k,0)\in[0,1)^n\times K$. Since the norm is the max
norm and $\delta\le1$, $\overline B(x',\delta)\subset C$, and $\overline B(x,\delta)=\overline B(x',\delta)+(k,0)\subset W$.
By periodicity $F(y+(k,0))=F(y)$ on $\overline B(x',\delta)$, so $\norm F\le M$ and $F$ is $L_C$-Lipschitz on $\overline B(x,\delta)$.
\end{proof}

\subsection{Local existence}\label{bpN:p3:s-exist}

\begin{lemma}[A priori short-time confinement]\label{bpN:p3:gen-short}
Let $J$ be an interval, $t_0\in J$, $\alpha:\R\to E$ with $\alpha(t_0)=x$ and $\alpha'(t)=F(\alpha(t))$ for all $t\in J$
(derivative two-sided; $F$ is only assumed to be a function), and suppose $\norm{F}\le M$ on $\overline B(x,\delta)$ with
$M>0$. Then $\norm{\alpha(t)-x}\le M\abs{t-t_0}$ for all $t\in J$ with $\abs{t-t_0}<\delta/M$; in particular
$\alpha(t)\in\overline B(x,\delta)$.
\source{Standard.}
\end{lemma}
\begin{proof}
\proofstep{Forward bootstrap}\label{step:p3:short-boot} Let $t\in J$, $t_0<t<t_0+\delta/M$ (the case $t<t_0$ follows by Lemma~\ref{bpN:p3:gen-symmetry}(b) and the case $t=t_0$ is trivial).
Let $A=\{s\in[t_0,t]:\norm{\alpha(u)-x}\le M(u-t_0)\ \forall u\in[t_0,s]\}$ and $s^*=\sup A$. $t_0\in A$ and $A$ is closed by
continuity of $\alpha$ on $J$, so $s^*\in A$.

\proofstep{Extension}\label{step:p3:short-ext} If $s^*<t$, then $\norm{\alpha(s^*)-x}\le M(s^*-t_0)<\delta$, so by continuity there is $\eta>0$, $s^*+\eta\le t$, with $\alpha([s^*,s^*+\eta])\subset\overline B(x,\delta)$. On this interval $\norm{\alpha'}=\norm{F(\alpha)}\le M$, so by
\ref{bpN:p3:ext-mvt} $\norm{\alpha(u)-\alpha(s^*)}\le M(u-s^*)$ and hence $\norm{\alpha(u)-x}\le M(u-t_0)$ there,
contradicting the definition of $s^*$. So $s^*=t$.
\end{proof}

\begin{proposition}[Local existence with uniform time]\label{bpN:p3:gen-uniform-exist}
Let $F$ be $C^1$ on the open set $W$ and $\Phi$ have uniform data $(\delta,M,L)$; put $M'=\max(M,1)$ and
$\varepsilon=\delta/M'$. Then for every $x\in\Phi$ and every $t_0\in\R$ there is $\alpha:\R\to E$ with $\alpha(t_0)=x$,
$\norm{\alpha(t)-x}<\delta$ and
$\texttt{IsSolIn}\,F\,W\,(t_0-\varepsilon,t_0+\varepsilon)\,\alpha$.
\source{Standard (Picard--Lindel\"of).}
\end{proposition}
\begin{proof}
\proofstep{Picard--Lindel\"of}\label{step:p3:ue-pl} Apply \ref{bpN:p3:ext-pl} with $a=\delta$, bound $M'$, Lipschitz constant $L$ and $T=\varepsilon$ ($M'T=\delta\le a$):
there is $\alpha$ with $\alpha(t_0)=x$ and derivative $F(\alpha(t))$ within $[t_0-\varepsilon,t_0+\varepsilon]$ at each point of it.
At points $t$ of the open interval this is a two-sided derivative, since the closed interval is a neighbourhood of $t$.

\proofstep{Staying in $W$}\label{step:p3:ue-stay} By Lemma~\ref{bpN:p3:gen-short} on $J=(t_0-\varepsilon,t_0+\varepsilon)$, $\norm{\alpha(t)-x}\le M'\abs{t-t_0}<\delta$, so $\alpha(t)\in\overline B(x,\delta)\subset W$.
\end{proof}

\begin{theorem}[Local existence, both time directions]\label{bpN:p3:gen-local}
Let $F$ be $C^1$ on the open set $W$, $x_0\in W$, $t_0\in\R$. There are $\varepsilon>0$ and $\alpha:\R\to E$ with $\alpha(t_0)=x_0$ and
$\texttt{IsSolIn}\,F\,W\,(t_0-\varepsilon,t_0+\varepsilon)\,\alpha$. Moreover $\varepsilon$ may be taken uniform on any compact
$K\subset W$ (for all $x_0\in K$).
\source{Standard.}
\end{theorem}
\begin{proof}
\proofstep{Reduction}\label{step:p3:loc-red} $\{x_0\}$ (resp.\ $K$) is compact; Lemma~\ref{bpN:p3:gen-uniform-compact} gives uniform data and Proposition~\ref{bpN:p3:gen-uniform-exist} the claim.
\end{proof}

\subsection{Uniqueness on every interval}\label{bpN:p3:s-unique}

\begin{theorem}[Uniqueness on arbitrary intervals]\label{bpN:p3:gen-uniqueness}
Let $F$ be $C^1$ on the open set $W$, $J\subset\R$ an interval (of any kind: open, closed, half-open, unbounded, or
a point), $t_0\in J$, and $y_1,y_2:\R\to E$ with $\texttt{IsSolIn}\,F\,W\,J\,y_i$ and $y_1(t_0)=y_2(t_0)$. Then
$y_1=y_2$ on $J$. No confinement, boundedness or smallness of either solution is assumed.
In particular, if $y_1,y_2$ are solutions on intervals $J_1,J_2$ and agree at one point of $J_1\cap J_2$, they agree on
$J_1\cap J_2$ (an interval), and two solutions which agree at one time of $J$ agree at all times of $J$.
\source{Standard; used for BM11 Theorem~2.1 (uniqueness) and in the maximal-solution argument of BM11 p.~13.}
\end{theorem}
\begin{proof}
\proofstep{A compact piece}\label{step:p3:uq-compact} Let $t_1\in J$, $t_1>t_0$ ($t_1<t_0$ is the mirror image, using the left-derivative form of \ref{bpN:p3:ext-unique}; $t_1=t_0$ is the hypothesis).
Then $[t_0,t_1]\subset J$, each $y_i$ is continuous on it, so $K=y_1([t_0,t_1])\cup y_2([t_0,t_1])$ is compact
(\ref{bpN:p3:ext-compact}) and $K\subset W$.

\proofstep{Lipschitz on $K$}\label{step:p3:uq-lip} Lemma~\ref{bpN:p3:gen-lip} gives $L$ with $F$ $L$-Lipschitz on $K$.

\proofstep{Apply \ref{bpN:p3:ext-unique}}\label{step:p3:uq-apply} Both $y_i$ are continuous on $[t_0,t_1]$, have right derivative $F(y_i(t))$ on $[t_0,t_1)$ (a two-sided derivative gives a right derivative) and take values in $K$; so
$y_1(t_1)=y_2(t_1)$.
\end{proof}

\subsection{Maximal solutions and continuation}\label{bpN:p3:s-max}

\begin{lemma}[Extension to an open interval]\label{bpN:p3:gen-extend}
Let $F$ be $C^1$ on the open set $W$, $J\ne\emptyset$ an interval and $\texttt{IsSolIn}\,F\,W\,J\,y$. Then there are an
open interval $J'\supset J$ and $y'$ with $y'=y$ on $J$ and $\texttt{IsSolIn}\,F\,W\,J'\,y'$.
\end{lemma}
\begin{proof}
\proofstep{Right end}\label{step:p3:ext-right} If $J$ has a largest element $b$, apply Theorem~\ref{bpN:p3:gen-local} at $x_0=y(b)\in W$, $t_0=b$: $\alpha$ on $(b-\varepsilon,b+\varepsilon)$.
Put $y_1(s)=y(s)$ for $s\le b$ and $y_1(s)=\alpha(s)$ for $s>b$. Then $y_1$ is a solution on $J_1=J\cup(b,b+\varepsilon)$: at $t\in J$, $t<b$, and at $t\in(b,b+\varepsilon)$ it agrees with $y$, resp.\ $\alpha$, near $t$;
at $t=b$ it has derivative $F(y(b))$ within $(-\infty,b]$ (it equals $y$ there) and within $[b,\infty)$ (it equals $\alpha$ there, as $\alpha(b)=y(b)$), so by \ref{bpN:p3:ext-calc} it has the two-sided derivative.
If $J$ has no largest element, $J_1=J$, $y_1=y$.

\proofstep{Left end}\label{step:p3:ext-left} Apply Step~\ref{step:p3:ext-right} to the reversed solution $y_1(-\cdot)$ (Lemma~\ref{bpN:p3:gen-symmetry}(b)). The result has no endpoint, so is an open interval.
\end{proof}

\begin{theorem}[Maximal solution]\label{bpN:p3:gen-maximal}
Let $F$ be $C^1$ on the open set $W$ and $x_0\in W$. There are $T_-\in[-\infty,0)$, $T_+\in(0,+\infty]$ and
$z:\R\to E$ with $z(0)=x_0$ and $\texttt{IsSolIn}\,F\,W\,(T_-,T_+)\,z$, such that every solution $y$ with
$y(0)=x_0$ on an interval $J\ni0$ satisfies $J\subset(T_-,T_+)$ and $y=z$ on $J$. Then $(T_-,T_+)$ is the
\emph{maximal interval} and $z|_{(T_-,T_+)}$ the \emph{maximal solution}; both are unique. For initial time $t_0$
instead of $0$ use Lemma~\ref{bpN:p3:gen-symmetry}(a).
\source{BM11 p.~13 (``maximal time of existence''); standard.}
\end{theorem}
\begin{proof}
\proofstep{The union}\label{step:p3:max-union} Let $\mathcal J$ be the set of open intervals $J\ni0$ carrying a solution $y_J$ with $y_J(0)=x_0$; $\mathcal J\ne\emptyset$ by Theorem~\ref{bpN:p3:gen-local}.
The union $I$ of $\mathcal J$ is an open interval $(T_-,T_+)\ni0$ (a union of intervals containing $0$).

\proofstep{Gluing}\label{step:p3:max-glue} For $t\in I$ choose $J\in\mathcal J$ with $t\in J$ and set $z(t)=y_J(t)$ (and $z(t)=x_0$ outside $I$). This is independent of $J$: for $J,J'\in\mathcal J$, $y_J=y_{J'}$ on $J\cap J'\ni0$ by Theorem~\ref{bpN:p3:gen-uniqueness}.
Each $t\in I$ has an open neighbourhood $J\in\mathcal J$ on which $z=y_J$, so $z$ is a solution on $I$ (local nature, (iii) above).

\proofstep{Maximality}\label{step:p3:max-max} If $y$ is a solution on an interval $J\ni0$ with $y(0)=x_0$, Lemma~\ref{bpN:p3:gen-extend} gives an open $J'\supset J$ and a solution $y'$ extending $y$, so $J'\in\mathcal J$, $J\subset J'\subset I$, and $y=y'=z$ on $J$ (Theorem~\ref{bpN:p3:gen-uniqueness}).
Uniqueness of $(T_-,T_+)$ and $z$ is immediate from the maximality property applied twice.
\end{proof}

\begin{theorem}[Continuation criterion, uniform form]\label{bpN:p3:gen-continuation}
Let $F$ be $C^1$ on the open set $W$, $x_0\in W$, $(T_-,T_+)$ and $z$ the maximal interval and solution of
Theorem~\ref{bpN:p3:gen-maximal}. Let $\Phi\subset E$ have uniform data $(\delta,M,L)$
(Definition~\ref{bpN:p3:gen-uniform}). If $T_+<\infty$ there is no $t_*\in[0,T_+)$ with $z(t)\in\Phi$ for all
$t\in[t_*,T_+)$. Symmetrically, if $T_->-\infty$ there is no $t_*\in(T_-,0]$ with $z(t)\in\Phi$ for all $t\in(T_-,t_*]$.
Equivalently: if the solution stays in $\Phi$ near the right (left) end of its maximal interval, that end is $+\infty$ ($-\infty$).
\source{Standard; BM11 p.~13 (``this solution is defined for all time''), P\"oschel 1993 p.~198 (``$T=T_*$'').}
\end{theorem}
\begin{proof}
\proofstep{Set-up}\label{step:p3:cont-setup} Suppose $T_+<\infty$ and $z(t)\in\Phi$ on $[t_*,T_+)$. Let $\varepsilon=\delta/\max(M,1)$ from Proposition~\ref{bpN:p3:gen-uniform-exist}. Choose $t_1\in[t_*,T_+)$ with $t_1>T_+-\varepsilon$.

\proofstep{A solution through $z(t_1)$}\label{step:p3:cont-alpha} Proposition~\ref{bpN:p3:gen-uniform-exist} with $x=z(t_1)\in\Phi$ and $t_0=t_1$ gives $\alpha$ on $(t_1-\varepsilon,t_1+\varepsilon)$ with $\alpha(t_1)=z(t_1)$.
By Theorem~\ref{bpN:p3:gen-uniqueness} on $(t_1-\varepsilon,t_1+\varepsilon)\cap(T_-,T_+)$, $\alpha=z$ there.

\proofstep{Gluing past $T_+$}\label{step:p3:cont-glue} Put $y=z$ on $(T_-,T_+)$ and $y=\alpha$ on $[T_+,t_1+\varepsilon)$; then $y=\alpha$ on $(t_1-\varepsilon,t_1+\varepsilon)$ and $y=z$ on $(T_-,T_+)$, two open sets covering $(T_-,t_1+\varepsilon)$ on each of which $y$ is a solution.
So $y$ is a solution on $(T_-,t_1+\varepsilon)\ni0$ with $y(0)=x_0$, and $t_1+\varepsilon>T_+$ contradicts Theorem~\ref{bpN:p3:gen-maximal}.
The left end follows by Lemma~\ref{bpN:p3:gen-symmetry}(b).
\end{proof}

\begin{corollary}[Compact continuation]\label{bpN:p3:gen-continuation-compact}
In the setting of Theorem~\ref{bpN:p3:gen-continuation}, if $K\subset W$ is compact and $z(t)\in K$ for all
$t\in[t_*,T_+)$ (some $t_*\in[0,T_+)$), then $T_+=+\infty$; likewise on the left.
\end{corollary}
\begin{proof}
Lemma~\ref{bpN:p3:gen-uniform-compact} gives uniform data for $K$; apply Theorem~\ref{bpN:p3:gen-continuation}.
\end{proof}

\begin{corollary}[Closed sets, periodic fields]\label{bpN:p3:gen-continuation-periodic}
Let $E=\R^n\times\R^n$, $U\subset\R^n$ open, $W=\R^n\times U$, $F$ $C^1$ on $W$ with $F(\theta+k,I)=F(\theta,I)$ for $k\in\Z^n$. If the maximal solution through $x_0\in W$ satisfies $z(t)_2\in K$ for $t\in[t_*,T_+)$ with $K\subset U$ compact, then $T_+=+\infty$ (likewise on the left). The angle components $z(t)_1$ may be unbounded.
\end{corollary}
\begin{proof}
Lemma~\ref{bpN:p3:gen-uniform-periodic} gives uniform data for $\Phi=\R^n\times K$; apply Theorem~\ref{bpN:p3:gen-continuation}.
\end{proof}

\begin{remark}[What the closed-set variant really needs]\label{bpN:p3:gen-uniform-remark}
Staying in a closed non-compact set $\Phi\subset W$ does \emph{not} suffice in general (the solution of $\dot x=x^2$
on $E=\R$, $W=\R$, $x_0=1$ stays in the closed set $[1,\infty)$ and blows up at $T_+=1$). What suffices is the
uniform data $(\delta,M,L)$ of Definition~\ref{bpN:p3:gen-uniform}: a $\delta$-neighbourhood of $\Phi$ inside $W$ on which $F$ is bounded and Lipschitz, uniformly. This holds for compact sets, for angle-periodic fields over a compact action set (Lemma~\ref{bpN:p3:gen-uniform-periodic}), and for the Newtonian regions bounded away from collisions with bounded velocities (Corollary~\ref{bpN:p3:newton-continuation}).
\end{remark}

\begin{theorem}[Confinement glue]\label{bpN:p3:gen-confinement}
Let $F$ be $C^1$ on the open set $W$, $x_0\in W$, $\tau>0$, and $\Phi\subset E$ a set with uniform data (e.g.\ compact
$\Phi\subset W$, or as in Lemma~\ref{bpN:p3:gen-uniform-periodic}). Assume the \emph{a priori confinement}: every
solution $y$ on an open interval $J\ni0$ with $y(0)=x_0$ satisfies $y(t)\in\Phi$ for all $t\in J$ with $\abs t\le\tau$.
Then the maximal interval contains $[-\tau,\tau]$ in its interior, so $z$ is a solution on $[-\tau,\tau]$ with
$z(t)\in\Phi$ for $\abs t\le\tau$, and every solution on $[-L,L]$ ($0\le L\le\tau$) with initial value $x_0$ equals $z$
there.
\source{BM11 p.~13 (the two-case argument ``Therefore this solution is defined for all time''); P\"oschel 1993 p.~198.}
\end{theorem}
\begin{proof}
\proofstep{No exit before $\tau$}\label{step:p3:conf-noexit} Suppose $T_+\le\tau$. The maximal solution $z$ on $(T_-,T_+)$ satisfies, for $t\in[0,T_+)$ ($\abs t<\tau$), $z(t)\in\Phi$ by confinement.
Theorem~\ref{bpN:p3:gen-continuation} gives $T_+=\infty$, a contradiction. So $T_+>\tau$, and likewise $T_->-\tau$.

\proofstep{Conclusion}\label{step:p3:conf-concl} $[-\tau,\tau]\subset(T_-,T_+)$ so $z$ is a solution on $[-\tau,\tau]$ and $z(t)\in\Phi$ for $\abs t\le\tau$ by confinement applied to $z$ itself. Uniqueness is Theorem~\ref{bpN:p3:gen-uniqueness}.
\end{proof}

\subsection{First exit from an open set; barriers}\label{bpN:p3:s-exit}

\begin{lemma}[First exit]\label{bpN:p3:gen-exit}
Let $\texttt{IsSolIn}\,F\,W\,J\,z$ with $J$ an interval, $0\in J$, $V\subset E$ open and $z(0)\in V$. Put
$A=\{t\in J:t\ge0,\ z(t)\notin V\}$. Then exactly one of the following holds.
\begin{enumerate}[label=(\roman*)]
\item $A=\emptyset$: $z(t)\in V$ for all $t\in J$, $t\ge0$.
\item $A\ne\emptyset$. Then $t^*=\min A$ exists, $t^*>0$, $t^*\in J$, $z(t)\in V$ for $t\in[0,t^*)$ and
  $z(t^*)\in\partial V=\overline V\setminus V$. The number $t^*$ is the \emph{first exit time} and equals
  $\sup\{t\in J:t\ge0,\ z([0,t])\subset V\}$; in particular $t^*$ is finite and lies in $J$ ($t^*$ is ``inside the
  solution's interval'' and the solution is defined, continuous, at $t^*$).
\end{enumerate}
The same holds backwards in time (replace $z(t)$ by $z(-t)$, Lemma~\ref{bpN:p3:gen-symmetry}(b)).
Only continuity of $z$ on $J$ is used; $F$ and $W$ play no role.
\source{Standard; used for BM11 p.~13 (``smallest time $t_*$ such that $\omega(t_*)\in R_K$'') and P\"oschel 1993 p.~198--199 (stopping-time arguments).}
\end{lemma}
\begin{proof}
\proofstep{Closedness}\label{step:p3:exit-closed} $z$ is continuous on $J$ (consequence (ii)), $E\setminus V$ is closed, so $A$ is closed in $J\cap[0,\infty)$.

\proofstep{The minimum}\label{step:p3:exit-min} If $A\ne\emptyset$, $t^*=\inf A\in[0,\infty)$; $t^*\in J$ since $J$ is an interval containing $0$ and a point of $A$ that is $\ge t^*$. By closedness $t^*\in A$, i.e.\ $z(t^*)\notin V$. As $z(0)\in V$, $t^*>0$.

\proofstep{Boundary point}\label{step:p3:exit-bdry} For $t\in[0,t^*)$, $t\notin A$, so $z(t)\in V$. By continuity $z(t^*)=\lim_{t\uparrow t^*}z(t)\in\overline V$, and $z(t^*)\notin V$; as $V$ is open, $\overline V\setminus V=\partial V$.

\proofstep{The supremum formula}\label{step:p3:exit-sup} $z([0,t])\subset V$ iff $t<t^*$ (for $t\in J$, $t\ge 0$), so the supremum is $t^*$.
\end{proof}

\begin{definition}[Invariant sets]\label{bpN:p3:gen-invariant}
$S\subset E$ is \emph{invariant} for $(F,W)$ if for every interval $J'$ and every $y$ with $\texttt{IsSolIn}\,F\,W\,J'\,y$:
if $y(t_1)\in S$ for one $t_1\in J'$ then $y(t)\in S$ for all $t\in J'$. (Equilibria $\{x:F(x)=0\}$ are invariant by
Theorem~\ref{bpN:p3:gen-uniqueness}; so is any union of such sets and any set $\{G=c\}$ for a first integral $G$.)
\end{definition}

\begin{corollary}[Barrier]\label{bpN:p3:gen-barrier}
Let $F$ be $C^1$ on the open set $W$, $V$ open, $\texttt{IsSolIn}\,F\,W\,J\,z$ with $J$ an interval, $0\in J$, $z(0)\in V$.
Let $G\subset E$ with $z(J)\subset G$, and suppose $\partial V\cap G\subset S$ where $S$ is invariant and $z(0)\notin S$
(e.g.\ $G$ an energy level, Theorem~\ref{bpN:p3:ham-energy}, and $\partial V\cap G$ a union of invariant sets not meeting the
initial point). Then $z(t)\in V$ for all $t\in J$ (both time directions).
\end{corollary}
\begin{proof}
\proofstep{Forward}\label{step:p3:barrier-fwd} If not, Lemma~\ref{bpN:p3:gen-exit}(ii) gives $t^*\in J$ with $z(t^*)\in\partial V\cap G\subset S$; invariance gives $z(0)\in S$, a contradiction.

\proofstep{Backward}\label{step:p3:barrier-bwd} Apply Step~\ref{step:p3:barrier-fwd} to $t\mapsto z(-t)$, a solution of $-F$ on $-J$; $G$, $\partial V$ are unchanged and $S$ is invariant for $-F$ as well (reversal maps solutions of $F$ to solutions of $-F$ and back).
\end{proof}

\subsection{Hamiltonian layer}\label{bpN:p3:s-ham}

Phase space $P=\R^n\times\R^n$, points $z=(\theta,I)$, max norm. $X_H$ is the frozen
\Lean{KAM.hamiltonianVectorField}, Definition~\ref{bpN:p0:vf} (convention of \Kref{bp:p0:hamiltonian}):
$X_H(z)=\bigl((\partial_{I_i}H(z))_i,(-\partial_{\theta_i}H(z))_i\bigr)$, with $\partial_{I_i}H(z)=DH(z)(0,e_i)$ and
$\partial_{\theta_i}H(z)=DH(z)(e_i,0)$.

\begin{lemma}[Regularity of $X_H$]\label{bpN:p3:ham-smooth}
Let $W\subset P$ be open and $H:P\to\R$ be $C^{k+1}$ on $W$, $k\ge0$. Then $X_H$ is $C^k$ on $W$; for $k\ge1$ it is $C^1$
(hence locally Lipschitz) on $W$, and $H$ is differentiable at every point of $W$.
\source{Standard.}
\end{lemma}
\begin{proof}
\proofstep{Derivative is $C^k$}\label{step:p3:hs-deriv} $DH:W\to\mathcal L(P,\R)$ is $C^k$ on the open $W$ (\ref{bpN:p3:ext-calc}, \Lean{ContDiffOn.fderiv\_of\_isOpen}).

\proofstep{Components}\label{step:p3:hs-comp} Each coordinate of $X_H$ is $\pm$ the composition of $DH$ with the continuous linear evaluation $\ell\mapsto\ell(0,e_i)$ or $\ell\mapsto\ell(e_i,0)$; hence $C^k$, and a map into $\R^n\times\R^n$ with $C^k$ coordinates is $C^k$.
\end{proof}

\begin{theorem}[Energy conservation]\label{bpN:p3:ham-energy}
Let $H:P\to\R$, $J\subset\R$ an interval and $z:\R\to P$ such that for every $t\in J$, $H$ is differentiable at $z(t)$
and $z$ has derivative $X_H(z(t))$ at $t$. Then $t\mapsto H(z(t))$ is constant on $J$.
(No regularity of $H$ beyond the stated pointwise differentiability, and no confinement.)
\source{BM11 p.~13 and P\"oschel 1993 p.~198 (``energy conservation'', ``$\Delta H=H(I(t),\theta(t))|_0^T$ vanishes''); standard.}
\end{theorem}
\begin{proof}
\proofstep{Chain rule}\label{step:p3:en-chain} For $t\in J$, $\frac{d}{dt}H(z(t))=DH(z(t))\bigl[X_H(z(t))\bigr]$ (\ref{bpN:p3:ext-calc}).

\proofstep{Cancellation}\label{step:p3:en-cancel} Write $z=z(t)$, $\ell=DH(z)$ and $X_H(z)=(a,b)$ with $a_i=\ell(0,e_i)$, $b_i=-\ell(e_i,0)$. By linearity, $\ell(a,b)=\sum_ia_i\,\ell(e_i,0)+\sum_ib_i\,\ell(0,e_i)=\sum_i\ell(0,e_i)\ell(e_i,0)-\sum_i\ell(e_i,0)\ell(0,e_i)=0$.

\proofstep{Constancy}\label{step:p3:en-const} Let $s<t$ in $J$. $H\circ z$ is continuous on $[s,t]\subset J$ and has derivative $0$ at every point, so $H(z(t))=H(z(s))$ (\ref{bpN:p3:ext-calc}).
\end{proof}

\begin{lemma}[Periodicity in the angles passes to solutions]\label{bpN:p3:ham-periodic}
Let $W\subset P$ be open with $W+(k,0)=W$ for all $k\in\Z^n$, and $H:P\to\R$ with $H(\theta+k,I)=H(\theta,I)$ for $(\theta,I)\in W$,
$k\in\Z^n$. Then $X_H(\theta+k,I)=X_H(\theta,I)$ on $W$. Consequently, if $\texttt{IsSolIn}\,X_H\,W\,J\,z$ then
$\texttt{IsSolIn}\,X_H\,W\,J\,(z+(k,0))$ for every $k\in\Z^n$, and $H\circ(z+(k,0))=H\circ z$.
\source{Standard.}
\end{lemma}
\begin{proof}
\proofstep{Field}\label{step:p3:per-field} Let $\tau_k(x)=x+(k,0)$. On the open set $W$, $H\circ\tau_k=H$, so $DH=D(H\circ\tau_k)=DH\circ\tau_k$ at each point of $W$ (the derivative of a translate; both sides are $0$ if $H$ is not differentiable, by the definition of \Lean{fderiv}). Evaluating on $(0,e_i)$ and $(e_i,0)$ gives $X_H\circ\tau_k=X_H$ on $W$.

\proofstep{Solutions}\label{step:p3:per-sol} Lemma~\ref{bpN:p3:gen-symmetry}(c) with $v=(k,0)$ applies to $F=X_H$.
\end{proof}

\subsection{The frozen solution predicates are instances}\label{bpN:p3:s-frozen}

\begin{proposition}[Track N and track I Hamiltonian predicates]\label{bpN:p3:ham-reduction}
Let $H:P\to\R$, $U\subset\R^n$ open, $W\subset P$ open with $\R^n\times U\subset W$, and $H$ be $C^2$ on $W$. Put
$W_U=\R^n\times U$. For every $J\subset\R$ and $z:\R\to P$,
\[
\texttt{IsSolutionOn}\,H\,U\,J\,z\iff\texttt{IsSolIn}\,X_H\,W_U\,J\,z
\iff\texttt{IsHamiltonianSolutionOn}\,H\,U\,J\,z,
\]
and $X_H$ is $C^1$ on $W_U$. (The three predicates are literally the same formula up to the letter naming the action set,
\Lean{KAM.Nekhoroshev.IsSolutionOn} and \Lean{KAM.Iso.IsHamiltonianSolutionOn}.)
\frozen{KAM.Nekhoroshev.IsSolutionOn, KAM.Iso.IsHamiltonianSolutionOn}
\source{Ours (bookkeeping).}
\end{proposition}
\begin{proof}
\proofstep{Membership}\label{step:p3:red-mem} $z(t)_2\in U\iff z(t)\in W_U$ by definition of $W_U$.

\proofstep{Differentiability is free}\label{step:p3:red-diff} For $z(t)\in W_U\subset W$, $H$ is $C^2$ on the open $W$, hence differentiable at $z(t)$ (Lemma~\ref{bpN:p3:ham-smooth}); so the clause $\texttt{DifferentiableAt}\,\R\,H\,(z\,t)$ is implied by $z(t)\in W_U$, and the two predicates coincide. Smoothness of $X_H$: Lemma~\ref{bpN:p3:ham-smooth}, $k=1$, restricted to $W_U\subset W$.
\end{proof}

\begin{proposition}[Newtonian predicates]\label{bpN:p3:newton-reduction}
Let $N\ge2$, $G\in\R$, $\texttt{mass}:\text{Fin}\,N\to\R$. Let $E_N=(\text{Fin}\,N\to\R^2)\times(\text{Fin}\,N\to\R^2)$ (max norm), $\ell=\texttt{planeLength}$,
\[
W_N=\{(q,v):\ell(q_j-q_i)>0\ \forall i\ne j\},\qquad
F_N(q,v)=\Bigl(v,\ \bigl(\textstyle\sum_{j\ne i}\tfrac{G\,\mathrm{mass}_j}{\ell(q_j-q_i)^3}(q_j-q_i)\bigr)_i\Bigr).
\]
Then $W_N$ is open, $F_N$ is $C^\infty$ on $W_N$, and for $Q,V:\R\to(\text{Fin}\,N\to\R^2)$ with $x(t)=(Q(t),V(t))$:
$\texttt{IsNewtonSolution}\,G\,\mathrm{mass}\,Q\,V\iff\texttt{IsSolIn}\,F_N\,W_N\,\R\,x$.
Likewise for \Lean{IsFixedSolution}\,$G\,m_0\,m\,\varepsilon$ with $W_N'=W_N\cap\{\ell(q_i)>0\ \forall i\}$ and the field
$F'(q,v)=(v,(-\frac{Gm_0}{\ell(q_i)^3}q_i+\varepsilon\sum_{j\ne i}\frac{G m_j}{\ell(q_j-q_i)^3}(q_j-q_i))_i)$, and for \Lean{IsRestrictedSolutionOn}\,$\mu\,J$ with
$W_\mu=\{(q,p):\texttt{CollisionFree}\,\mu\,(q,p)\}$ and the total extension $F_\mu$ of \Lean{restrictedVectorField}\,$\mu$ by the same formula
(with the values off $W_\mu$ irrelevant): $\texttt{IsRestrictedSolutionOn}\,\mu\,J\,z\iff\texttt{IsSolIn}\,F_\mu\,W_\mu\,J\,z$.
In all three cases $F$ is $C^\infty$ on the open set $W$ and the general layer applies with $J=\R$ (resp.\ the given $J$).
\frozen{KAM.Planetary.IsNewtonSolution, KAM.Planetary.IsFixedSolution, KAM.Iso.IsRestrictedSolutionOn}
\source{Ours (bookkeeping); Fejoz Section~6.1, Newton equations as in the frozen text.}
\end{proposition}
\begin{proof}
\proofstep{Open domain and smoothness}\label{step:p3:nr-smooth} $\ell(w)^2=\texttt{planeSq}(w)=w_0^2+w_1^2$ is a polynomial, so $\{\ell>0\}=\{\texttt{planeSq}>0\}$ is open, and $\ell=\sqrt{\cdot}$ is $C^\infty$ where $\texttt{planeSq}>0$. Hence $\ell(q_j-q_i)^{-3}$ is $C^\infty$ on $W_N$ and $F_N$ is a finite sum of products of $C^\infty$ maps. Same for $F'$ and $F_\mu$ ($\texttt{radius}=\sqrt{\texttt{radiusSq}}$ with the shifts by $\texttt{primary}_{1,2}\mu$).

\proofstep{Derivatives of pairs}\label{step:p3:nr-pair} A map into a finite product has derivative $(a,b)$ at $t$ iff each component has the corresponding derivative (\Lean{hasDerivAt\_pi}, \Lean{HasDerivAt.prodMk} and converse). The first block of the frozen definition (distance positivity) is membership in $W_N$; the second block says $Q$ has derivative $V(t)$; the third says $V$ has derivative the force. Together they say exactly that $x$ has derivative $F_N(x(t))$ at $t$.
\end{proof}

\begin{corollary}[Newtonian continuation]\label{bpN:p3:newton-continuation}
In the setting of Proposition~\ref{bpN:p3:newton-reduction}, let $x$ be the maximal solution through $x_0\in W_N$
(Theorem~\ref{bpN:p3:gen-maximal}), $\rho>0$, $\bar V\ge0$ and
$\Phi_{\rho,\bar V}=\{(q,v):\ell(q_j-q_i)\ge\rho\ (i\ne j),\ \abs v\le\bar V\}$. If $T_+<\infty$ then for every $\rho,\bar V$
and every $t_*\in[0,T_+)$ there is $t\in[t_*,T_+)$ with $x(t)\notin\Phi_{\rho,\bar V}$ (a finite-time singularity needs
near-collisions or unbounded velocities; positions may be unbounded). The same for $F'$ with the additional condition $\ell(q_i)\ge\rho$, and for $F_\mu$ with $\ell\ge\rho$ from both primaries and $\abs q,\abs p\le\bar V$.
\end{corollary}
\begin{proof}
\proofstep{Uniform data}\label{step:p3:nc-data} Let $\delta=\rho/8$. If $\abs{y-x}\le\delta$ (max norm) with $x\in\Phi_{\rho,\bar V}$ then each $\ell(q'_j-q'_i)\ge\rho-2\sqrt2\,\delta\ge\rho/2$, so $y\in W_N$ and $\overline B(x,\delta)\subset W_N$.
On this ball $\ell\ge\rho/2$ and $\abs{v'}\le\bar V+\delta$; the entries of $F_N$ and $DF_N$ are polynomials in $v'$, the differences $q'_j-q'_i$ (Euclidean length $\le\ell(q'_j-q'_i)$) and the inverse powers $\ell^{-3},\ell^{-5}$ ($D(w/\ell^3)=\ell^{-3}I-3\ell^{-5}ww^{\top}$), hence bounded by a constant $C(N,G,\mathrm{mass},\rho,\bar V)$ \emph{independent of the position of $x$}. By \ref{bpN:p3:ext-mvt} on the convex ball, $F_N$ is $C$-Lipschitz there. So $\Phi_{\rho,\bar V}$ has uniform data.

\proofstep{Apply}\label{step:p3:nc-apply} Theorem~\ref{bpN:p3:gen-continuation} gives the claim; the other two fields are identical in structure (for $F_\mu$ the extra terms are linear in $(q,p)$ and bounded on $\abs q,\abs p\le\bar V$).
\end{proof}

\subsection{The track-N statements}\label{bpN:p3:s-N}

Throughout, $R,s>0$, $h:\C^n\to\C$ with $\texttt{IsAnalyticIntegrable}\,R\,s\,h$, $\varepsilon\ge0$ and $f:\C^n\times\C^n\to\C$ with
$\texttt{IsAnalyticPerturbation}\,R\,s\,\varepsilon\,f$ (Definitions~\ref{bpN:p0:integrable}, \ref{bpN:p0:pert}), $H=H_{h,f}$
(Definition~\ref{bpN:p0:realH}), $W_s=\R^n\times B(0,R+s)\subset P$ (open; $B(0,R+s)=\mathcal B_s(R)$ by Remark~\ref{bpN:p0:rdomain-eq}), and
$U\subset B(0,R+s)$ open (in N-HEAD-C, $U=B(0,R)$).

\begin{lemma}[Input from Part 1: $H$ is smooth on $W_s$]\label{bpN:p3:holo-input}
$H_{h,f}$ is $C^\infty$ (in particular $C^2$) on the open set $W_s$, $X_{H}=X_{h+f}$ on $W_s\cap(\R^n\times\R^n)$, and
$H(\theta+k,I)=H(\theta,I)$ for $(\theta,I)\in W_s$, $k\in\Z^n$.
\source{Part 1 forward label \Nref{bpN:p1:holo-real}; the two ingredients are \Kref{bp:p1:holomorphic-regularity}, \Kref{bp:p1:hamiltonian-real}; periodicity is Definition~\ref{bpN:p0:pert}(ii).}
\end{lemma}
\begin{proof}
\proofstep{Regularity}\label{step:p3:hi-reg} The function $G(\theta,I)=h(I)+f(\theta,I)$ is holomorphic on the open set $D_s(R)\supset W_s$ (real points $(\theta,I)$ with $I\in\mathcal B_s(R)=B(0,R)_s\cap\R^n$ lie in $T_s\times B(0,R)_s$), hence $C^\infty$ in the real sense (\Kref{bp:p1:holomorphic-regularity}); it is real on $W_s$ by Definition~\ref{bpN:p0:integrable}(iv) and \ref{bpN:p0:pert}(iii). So $H=\operatorname{Re}G=G$ on $W_s$ is the restriction of $C^\infty$ map to a real subspace: $C^\infty$ on $W_s$; $X_H=X_G$ at real points is \Kref{bp:p1:hamiltonian-real}. (Packaged as \Nref{bpN:p1:holo-real}.)

\proofstep{Periodicity}\label{step:p3:hi-per} $f(\theta+k,I)=f(\theta,I)$ for real $\theta$, $I\in\mathcal B_s(R)$, and $h$ depends on $I$ only.
\end{proof}

\begin{theorem}[$X_H$ is $C^1$, hence locally Lipschitz]\label{bpN:p3:lipschitz}
$X_{H_{h,f}}$ is $C^1$ on $W_s\supset\R^n\times B(0,R)$; consequently it is Lipschitz on every compact subset of $W_s$, and
on every set $\R^n\times K$ with $K\subset B(0,R+s)$ compact there are $\delta,M,L$ as in Lemma~\ref{bpN:p3:gen-uniform-periodic}.
\source{BM11 p.~13 (implicit in ``the solution''); standard.}
\end{theorem}
\begin{proof}
\proofstep{Smoothness}\label{step:p3:lip-smooth} Lemma~\ref{bpN:p3:holo-input} and Lemma~\ref{bpN:p3:ham-smooth} ($k=1$).
\proofstep{Lipschitz}\label{step:p3:lip-lip} Lemma~\ref{bpN:p3:gen-lip} on compact sets, Lemma~\ref{bpN:p3:gen-uniform-periodic} (with $U=B(0,R+s)$, using the periodicity of $X_H$ from Lemma~\ref{bpN:p3:ham-periodic} and Lemma~\ref{bpN:p3:holo-input}) for the cylinders.
\end{proof}

\begin{theorem}[Local existence and uniqueness for $\texttt{IsSolutionOn}$]\label{bpN:p3:local-exist}
For every $t_0\in\R$ and $z_0\in\R^n\times U$ there are $\epsilon>0$ and $y:\R\to P$ with $y(t_0)=z_0$ and
$\texttt{IsSolutionOn}\,H\,U\,(t_0-\epsilon,t_0+\epsilon)\,y$ (both time directions). If $J$ is any interval, $t_0\in J$ and
$y_1,y_2$ satisfy $\texttt{IsSolutionOn}\,H\,U\,J\,y_i$ with $y_1(t_0)=y_2(t_0)$, then $y_1=y_2$ on $J$.
\source{BM11 p.~13 and Theorem~2.1 p.~5; Picard--Lindel\"of (external fact \ref{bpN:p3:ext-pl}).}
\end{theorem}
\begin{proof}
\proofstep{Reduction}\label{step:p3:le-red} Proposition~\ref{bpN:p3:ham-reduction} (with $W=W_s$, $H$ $C^2$ there by Lemma~\ref{bpN:p3:holo-input}) turns $\texttt{IsSolutionOn}\,H\,U$ into $\texttt{IsSolIn}\,X_H\,(\R^n\times U)$ with $X_H$ $C^1$ on $\R^n\times U$.
\proofstep{Existence, uniqueness}\label{step:p3:le-apply} Theorem~\ref{bpN:p3:gen-local} and Theorem~\ref{bpN:p3:gen-uniqueness}.
\end{proof}

\begin{theorem}[Energy conservation for $\texttt{IsSolutionOn}$]\label{bpN:p3:energy}
If $\texttt{IsSolutionOn}\,H\,U\,J\,z$ for an interval $J$ (here $H$ may be any function $P\to\R$ and $U$ any set), then $H(z(t))=H(z(t'))$ for $t,t'\in J$.
In particular, for $H=H_{h,f}$ and $J=[0,T]$ the quantity $\Delta H=H(z(t))|_0^T$ of P\"oschel's Lemma~1 vanishes.
\source{P\"oschel 1993 p.~198 (``by energy conservation \dots $\Delta H$ vanishes''; page read); BM11 p.~13.}
\end{theorem}
\begin{proof}
\proofstep{Direct}\label{step:p3:en-direct} The definition gives differentiability of $H$ at every $z(t)$ and $z'(t)=X_H(z(t))$; Theorem~\ref{bpN:p3:ham-energy}.
\end{proof}

\begin{theorem}[Maximal solutions and continuation]\label{bpN:p3:continuation}
Let $z_0=(\theta_0,I_0)$ with $I_0\in U$. There is a unique maximal interval $(T_-,T_+)\ni0$ and a solution $z$ with
$z(0)=z_0$ and $\texttt{IsSolutionOn}\,H\,U\,(T_-,T_+)\,z$, such that every $y$ with $y(0)=z_0$ and $\texttt{IsSolutionOn}\,H\,U\,J\,y$
on an interval $J\ni0$ has $J\subset(T_-,T_+)$ and $y=z$ on $J$ (Theorem~\ref{bpN:p3:gen-maximal} through Proposition~\ref{bpN:p3:ham-reduction}).
If $K\subset U$ is compact and $z(t)_2\in K$ for all $t\in[t_*,T_+)$ for some $t_*\in[0,T_+)$, then $T_+=+\infty$; likewise on the left.
The angle components $z(t)_1$ may be unbounded: no obstruction arises because $X_H$ is $1$-periodic in $\theta$.
\source{BM11 p.~13 (``From this estimate we can deduce that the solution $I(t)$ belongs to some compact ball \dots this solution is defined for all time''); P\"oschel 1993 p.~198 (``$T=T_*$'').}
\end{theorem}
\begin{proof}
\proofstep{General statements}\label{step:p3:co-general} By Proposition~\ref{bpN:p3:ham-reduction}, $\texttt{IsSolutionOn}\,H\,U\,J\,y\iff\texttt{IsSolIn}\,X_H\,(\R^n\times U)\,J\,y$ and $X_H$ is $C^1$ on $W=\R^n\times U$; maximality is Theorem~\ref{bpN:p3:gen-maximal}.
\proofstep{Periodic cylinder}\label{step:p3:co-periodic} $X_H$ is $1$-periodic in $\theta$ on $W$ (Lemmas~\ref{bpN:p3:ham-periodic}, \ref{bpN:p3:holo-input}). Corollary~\ref{bpN:p3:gen-continuation-periodic} yields $T_+=\infty$.
\end{proof}

\begin{theorem}[Uniqueness on closed intervals]\label{bpN:p3:uniqueness-closed}
For $L\ge0$ and $y_1,y_2$ with $y_1(0)=y_2(0)$ and $\texttt{IsSolutionOn}\,H\,U\,[-L,L]\,y_i$ one has $y_1=y_2$ on $[-L,L]$. In particular a solution on $[-L,L]$ agrees with every solution on a larger interval containing $[-L,L]$ and the same initial value (clause (d) of Theorem~\ref{bpN:p0:head}, where $U=B(0,R)$, $L\in[0,\tau]$, $y$ arbitrary and $z$ the exhibited solution).
No confinement of $y$ is assumed: Gronwall is replaced by Theorem~\ref{bpN:p3:gen-uniqueness}, whose Lipschitz constant comes from the compact set $y_1([-L,L])\cup y_2([-L,L])$.
\source{Standard (Gronwall); clause (d) of the frozen N-HEAD-C.}
\end{theorem}
\begin{proof}
\proofstep{Apply}\label{step:p3:uc-apply} $J=[-L,L]$, $t_0=0$ in Theorem~\ref{bpN:p3:local-exist} (uniqueness part).
\end{proof}

\begin{theorem}[Confinement glue]\label{bpN:p3:confinement-glue}
Let $\theta_0\in\R^n$, $I_0\in U$, $\tau>0$ and $\rho>0$ with $\overline B(I_0,\rho)\subset U$ (e.g.\ $\abs{I_0}<R/2$, $\rho\le R/4$, $U=B(0,R)$).
Assume that for every open interval $J\ni0$ and every $y$ with $y(0)=(\theta_0,I_0)$ and $\texttt{IsSolutionOn}\,H\,U\,J\,y$ one has
$\abs{y(t)_2-I_0}\le\rho$ for all $t\in J$ with $\abs t\le\tau$. Then there is $z$ with $z(0)=(\theta_0,I_0)$,
$\texttt{IsSolutionOn}\,H\,U\,[-\tau,\tau]\,z$ and $\abs{z(t)_2-I_0}\le\rho$ for $\abs t\le\tau$; and for every $L\in[0,\tau]$ and every
$y$ with $y(0)=z(0)$ and $\texttt{IsSolutionOn}\,H\,U\,[-L,L]\,y$, $y=z$ on $[-L,L]$.
This is the assembly of clauses (a)--(d) of Theorem~\ref{bpN:p0:head} from an a priori estimate (Part~10); $\rho=C\varepsilon^{1/(2n)}$ there.
\source{BM11 p.~13 (two cases, ``Therefore this solution is defined for all time''); P\"oschel 1993 p.~198.}
\end{theorem}
\begin{proof}
\proofstep{Compact confinement set}\label{step:p3:cg-set} $K=\overline B(I_0,\rho)$ is compact and $\subset U$; $\Phi=\R^n\times K$ has uniform data by Lemma~\ref{bpN:p3:gen-uniform-periodic} (with $W=\R^n\times U$, $F=X_H$ periodic by Lemma~\ref{bpN:p3:ham-periodic}).
\proofstep{General glue}\label{step:p3:cg-glue} By Proposition~\ref{bpN:p3:ham-reduction} the hypothesis reads: every general solution on an open $J\ni0$ with this initial value stays in $\Phi$ for $\abs t\le\tau$.
Theorem~\ref{bpN:p3:gen-confinement} gives $z$ on $[-\tau,\tau]$ in $\Phi$, i.e.\ $\abs{z(t)_2-I_0}\le\rho$, and uniqueness on $[-L,L]$ (it is also Theorem~\ref{bpN:p3:uniqueness-closed}).
\end{proof}

\subsection{Label index for tracks I and P}\label{bpN:p3:s-index}

\begin{center}\small
\begin{tabular}{@{}p{0.4\textwidth}p{0.55\textwidth}@{}}
\texttt{bpN:p3:gen-sol}, \texttt{gen-symmetry} & predicate; time shift, reversal, space translation\\
\texttt{gen-lip}, \texttt{gen-uniform}, \texttt{gen-uniform-compact}, \texttt{gen-uniform-periodic} & Lipschitz and uniform data\\
\texttt{gen-short}, \texttt{gen-uniform-exist}, \texttt{gen-local} & local existence, uniform time\\
\texttt{gen-uniqueness} & uniqueness on all intervals, no confinement\\
\texttt{gen-extend}, \texttt{gen-maximal} & maximal solutions\\
\texttt{gen-continuation}, \texttt{-compact}, \texttt{-periodic}, \texttt{gen-confinement} & continuation, closed sets with uniform bounds, confinement glue\\
\texttt{gen-exit}, \texttt{gen-invariant}, \texttt{gen-barrier} & first exit, barrier\\
\texttt{ham-smooth}, \texttt{ham-energy}, \texttt{ham-periodic} & Hamiltonian layer\\
\texttt{ham-reduction}, \texttt{newton-reduction}, \texttt{newton-continuation} & frozen predicates of I, N, P\\
\texttt{holo-input}, \texttt{lipschitz}, \texttt{local-exist}, \texttt{energy}, \texttt{continuation}, \texttt{uniqueness-closed}, \texttt{confinement-glue} & track N specializations\\
\end{tabular}
\end{center}
